\documentclass[11pt]{article}
\usepackage[a4paper,margin=26mm]{geometry}
\usepackage{amsmath,amssymb,mathtools}
\usepackage{unicode-math}
\setmainfont{Latin Modern Roman}
\setmathfont{Latin Modern Math}
\usepackage{xurl,hyperref}
\hypersetup{hidelinks,pdftitle={UnsolvedMath Problem 10300031}}
\newcommand{\sref}[2]{\hyperlink{source:#1:#2}{[S#2]}}
\newcommand{\cataloguescope}[1]{#1}
\begin{document}
% UnsolvedMath ID 10300031; source catalogue code AMR-102-0031
\setcounter{section}{10300030}
\section{Leaf spaces and transverse structures — Question 8.8}\label{10300031}
{\small\sffamily UnsolvedMath ID 10300031 · Topology}
\subsection{Definitions and mathematical statement}

\paragraph{Shared notation.}$\mathbb N=\{1,2,\ldots\}$, and $\mathbb Z,\mathbb Q,\mathbb R,\mathbb C$ denote the integers, rationals, reals and complex numbers. Any other conventions are specified locally.


A lamination of $S^1$ is a closed family of unordered pairs of distinct circle points, represented by disjoint hyperbolic geodesics in the disk. Complementary regions have finite hyperbolic area. A transverse pair $\Lambda^+,\Lambda^-$ has transverse geodesic intersections between the two families. A group preserving both acts minimally on each family when every leaf orbit is dense in that leaf space. Two groups are commensurable when they have isomorphic finite-index subgroups.
\paragraph{Statement recorded in the supplied source.}
Given such a transverse pair of circle laminations with finite-area complementary regions, and a group $\Gamma\le\operatorname{Homeo}(S^1)$ preserving them and acting minimally on the leaves of each, characterize when $\Gamma$ is commensurable with $\pi_1(M)$ for a hyperbolic $3$-manifold $M$. \sref{10300031}{2}
\subsection{Short English statement}
Recognize when a group acting minimally on a pair of transverse circle laminations is commensurable with a hyperbolic three-manifold group.
\subsection{Sources}
\begin{description}
\item[\hypertarget{source:10300031:1}{S1}] ULAM.AI and the UnsolvedMath Contributors, \emph{UnsolvedMath: A Curated Collection of Open Mathematics Problems}, record 10300031 (AMR-102-0031). CC BY 4.0. \url{https://huggingface.co/datasets/ulamai/UnsolvedMath}.
\item[\hypertarget{source:10300031:2}{S2}] Danny Calegari, \emph{Problems in foliations and laminations of 3-manifolds}, version 0.78 (2002), Question 8.8; Definitions 1.1--1.10 and local definitions and remarks. \url{https://arxiv.org/abs/math/0209081v1}.
\end{description}
\label{end:10300031}
\end{document}
